\section{Conclusiones}

En el presente trabajo se implementó el ensayo \textit{Fiber-Bundle Pull-Out} (FBPO) como técnica
de caracterización interfacial a mesoescala en materiales compuestos reforzados con fibras,
cumpliendo el objetivo general de determinar la resistencia al corte interfacial (IFSS) y evaluar
la utilidad del método como herramienta confiable para el estudio de la adhesión fibra - matriz.

Se diseñó y fabricó un sistema experimental que incluyó mordazas, moldes y probetas, incorporando
criterios de precisión geométrica, reproducibilidad y compatibilidad con equipos de tracción
convencionales. El sistema desarrollado permitió la fabricación consistente de probetas y la
ejecución controlada de los ensayos FBPO, constituyendo una base experimental robusta para la
obtención de resultados repetibles y trazables.

El ensayo FBPO fue implementado exitosamente en probetas de material compuesto fabricadas con dos
matrices poliméricas y dos tipos de fibra de carbono. A partir de un total de 100 ensayos válidos,
se obtuvieron valores experimentales de IFSS aparente que describen la interacción fibra - matriz
y evidencian diferencias claras entre los sistemas evaluados, tanto en términos de magnitud como
de dispersión experimental. Además, las observaciones microscópicas post ensayo realizadas mediante
SEM permitieron identificar zonas donde la resina permanecía adherida, zonas de fibra libre de
matriz y regiones donde aún permanecía el teflón utilizado durante la fabricación, confirmando
la heterogeneidad de la interacción interfacial.

Los análisis estadísticos basados en ANOVA de Welch y comparaciones múltiples del tipo Games-Howell
indicaron que el tipo de fibra constituye el factor dominante en la respuesta interfacial dentro
del rango experimental estudiado, mientras que el efecto del tipo de matriz resulta secundario
cuando se emplea fibra genérica. La comparación entre los sistemas MEPOX-HTA y EPO-HTA mostró que
las diferencias en la adhesión dependen tanto de la matriz como de la combinación con la fibra de
grado aeroespacial, siendo observables a distintos niveles de aumento: de la hebra completa, por
zonas y a un solo filamento, lo que permitió una evaluación microscópica integral de la interfase.

La comparación de los valores obtenidos con rangos de IFSS reportados en la literatura para ensayos
FBPO permitió contextualizar los resultados experimentales y evidenció una concordancia razonable 
con estudios previos, reforzando la aplicabilidad del método implementado como técnica de
caracterización interfacial a mesoescala. Los hallazgos SEM aportan evidencia directa de los
distintos mecanismos de despegue y deslizamiento interfacial, complementando la información
obtenida a nivel macroscópico y reforzando la interpretación de los valores de IFSS.

En conjunto, los resultados obtenidos confirman que el ensayo FBPO, acompañado de un diseño
experimental cuidadosamente controlado y de un análisis estadístico apropiado, constituye una
herramienta válida y reproducible para la evaluación comparativa de la adhesión fibra - matriz
en materiales compuestos. Esta metodología permite sistematizar la caracterización interfacial
a mesoescala y proporciona un marco sólido para la interpretación de los mecanismos de interacción
observados, integrando tanto la evaluación macroscópica como la microscópica.

\clearpage
\subsection{Trabajo futuro}

A partir de los resultados obtenidos y de las limitaciones identificadas en el presente estudio,
se proponen diversas líneas de trabajo futuro orientadas a mejorar la robustez del método
experimental Fiber-Bundle Pull-Out (FBPO) y profundizar en la comprensión del comportamiento
interfacial.

En primer lugar, se plantea la optimización del diseño experimental, incluyendo la fabricación
de mordazas específicas para ensayos FBPO que mejoren la alineación de la hebra de fibras, reduzcan
efectos de flexión y minimicen daños durante la sujeción. Este refinamiento contribuiría a
disminuir la variabilidad experimental y aumentar la reproducibilidad de los resultados,
especialmente para sistemas con fibra genérica o mayor dispersión interfacial.

De manera complementaria, se propone explorar modelos numéricos orientados a la des- cripción de
la transferencia de carga y del despegue interfacial a mesoescala. Estos modelos permitirían
profundizar en la comprensión de los mecanismos físicos que gobiernan la interac- ción fibra - matriz
y complementar el análisis experimental, integrando los resultados de las micrografías con la
predicción de tensiones locales y comportamiento de la IFSS aparente.

En conjunto, estas líneas de trabajo futuro consolidarían el uso del ensayo FBPO como herramienta
reproducible y comparativa, contribuyendo a la sistematización del análisis de la adhesión
fibra - matriz y facilitando la extensión del método a diferentes configuraciones y condiciones
de procesamiento.
